% Artikel prosiding IEEE yang diringkas dari naskah tugas akhir.
% Kompilasi dari akar proyek:
%   latexmk -pdf -interaction=nonstopmode proceedings/main_proceedings.tex

\documentclass[conference]{IEEEtran}
\IEEEoverridecommandlockouts

\usepackage[T1]{fontenc}
\usepackage[utf8]{inputenc}
\usepackage{cite}
\usepackage{amsmath,amssymb,amsfonts}
\usepackage{graphicx}
\usepackage{textcomp}
\usepackage{booktabs}
\usepackage{array}
\usepackage{tabularx}
\usepackage{placeins}
\usepackage{float}
\usepackage{url}
\usepackage[hidelinks]{hyperref}

\graphicspath{{gambar/}{../gambar/}}

% Mendukung kompilasi dari akar proyek maupun langsung dari proceedings/.
\IfFileExists{proceedings/sections/00-abstract.tex}
  {\newcommand{\sectionspath}{proceedings/sections/}}
  {\newcommand{\sectionspath}{sections/}}

% Gunakan metadata tesis jika tersedia; fallback menjaga berkas proceedings
% tetap dapat dikompilasi secara mandiri.
\IfFileExists{pustaka/variables.tex}
  {\input{pustaka/variables.tex}}
  {\IfFileExists{../pustaka/variables.tex}
    {\input{../pustaka/variables.tex}}
    {\newcommand{\authorname}{Faiz Setya Kurniawan}
     \newcommand{\nrp}{5022221133}
     \newcommand{\advisor}{Astria Nur Irfansyah, S.T., M.Eng., Ph.D.}
     \newcommand{\advisornip}{198103252010121002}
     \newcommand{\coadvisor}{Dr. Ir. Hendra Kusuma, M.Eng.Sc.}
     \newcommand{\coadvisornip}{196409021989031003}}}

\providecommand{\citep}[1]{\cite{#1}}
\providecommand{\citet}[1]{\cite{#1}}
\providecommand{\parencite}[1]{\cite{#1}}

\renewcommand{\abstractname}{Abstrak}
\renewcommand{\IEEEkeywordsname}{Kata Kunci}
\renewcommand{\figurename}{Gambar}
\renewcommand{\tablename}{Tabel}
\renewcommand{\refname}{Daftar Pustaka}

% Header prosiding tiga baris pada setiap halaman. Jumlah \headheight dan
% \headsep dipertahankan 30 pt agar area teks IEEE tidak berubah.
\setlength{\headheight}{30pt}
\setlength{\headsep}{0pt}
\newcommand{\proceedingsheader}{%
  \parbox[b]{\textwidth}{%
    \raggedright\fontsize{10}{10}\selectfont
    \textbf{PROSIDING SEMINAR TUGAS AKHIR 134 (2026)}\\
    Departemen Teknik Elektro\\
    HIMATEKTRO ITS}}
\makeatletter
\def\ps@proceedings{%
  \def\@oddhead{\proceedingsheader}%
  \let\@evenhead\@oddhead
  \def\@oddfoot{}%
  \let\@evenfoot\@oddfoot}
\makeatother

\begin{document}

\pagestyle{proceedings}

\title{DESAIN BIQUAD LOW-PASS FILTER DAYA RENDAH DENGAN TEKNOLOGI CMOS OTA UNTUK REDUKSI NOISE PADA SINYAL ELEKTROKARDIOGRAM}

\author{%
  \makebox[\textwidth][c]{%
    \parbox[t]{0.31\textwidth}{%
      \centering
      \textit{Penulis}\\[0.25em]
      \parbox[t][2.5em][t]{\linewidth}{\centering\authorname}\\
      NRP \nrp\\
      Departemen Teknik Elektro\\
      Institut Teknologi Sepuluh Nopember\\
      Surabaya
    }\hfill
    \parbox[t]{0.31\textwidth}{%
      \centering
      \textit{Pembimbing}\\[0.25em]
      \parbox[t][2.5em][t]{\linewidth}{\centering\advisor}\\
      NIP \advisornip\\
      Departemen Teknik Elektro\\
      Institut Teknologi Sepuluh Nopember\\
      Surabaya
    }\hfill
    \parbox[t]{0.31\textwidth}{%
      \centering
      \textit{Ko-pembimbing}\\[0.25em]
      \parbox[t][2.5em][t]{\linewidth}{\centering\coadvisor}\\
      NIP \coadvisornip\\
      Departemen Teknik Elektro\\
      Institut Teknologi Sepuluh Nopember\\
      Surabaya
    }%
  }%
}

\maketitle
\thispagestyle{proceedings}

\input{\sectionspath 00-abstract.tex}
\input{\sectionspath 01-pendahuluan.tex}
\input{\sectionspath 02-metodologi.tex}
\input{\sectionspath 03-hasil-pembahasan.tex}
\input{\sectionspath 04-kesimpulan.tex}

\bibliographystyle{IEEEtran}
\IfFileExists{pustaka/pustaka.bib}
  {\bibliography{pustaka/pustaka}}
  {\bibliography{../pustaka/pustaka}}

\end{document}
